\section{Montesquieu on the betrayal of truth}

From Montesquieu's \textit{Spirit of the Laws}

\begin{quotex}
Dans les monarchies extrêmement absolues, les historiens trahissent la vérité, parce qu'ils n'ont pas la liberté de la dire : dans les Etats extrêmement libres, ils trahissent la vérité à cause de leur liberté même, qui, produisant toujours des divisions, chacun devient aussi esclave des préjugés de sa faction, qu'il le serait d'un despote.
\end{quotex}


\begin{quotex}
In absolute monarchies, historians betray the truth because they lack the freedom to say it. In politically free countries, they betray the truth because of their very freedom which then produces factions so that each one becomes a slave of his faction, which, for him, is equivalent to a despot.
\end{quotex}


I suppose it is pointless to comment since Montesquieu says it as well as it can be said. It is as useless to implore partisans to reject their partisanship as it would be to command the waves to recede. Partisans are under the sway of \emph{thymos}, which is the blind will to ``dominate''; philosophers are dominated by \emph{nous}, the pursuit of wisdom.


\flrightit{Posted on 2008-04-12 by Cologero}

